\documentclass[12pt,a4paper]{article}
\usepackage[utf8]{inputenc}
\usepackage[T1]{fontenc}
\usepackage{geometry}
\geometry{margin=2.5cm}

\begin{document}

\section*{4. Pick and Place}

The strategy for the pick and place task is built around two fixed Cartesian waypoints:
one above the pick-up location and one above the place location, both separated by a fixed
base-joint rotation of $2 \times 0.7$\,rad.  Using inverse kinematics these Cartesian poses
are converted to joint-space waypoints at initialisation time.  The robot then executes the
following sequence:

\begin{itemize}
    \item Move in joint space to above the pick-up location with the gripper open.

    \item Move straight down in Cartesian space until a collision is detected.

    \item Close the gripper until resistance is sensed.

    \item Move in joint space to the midpoint waypoint.

    \item Move in joint space to above the place location.

    \item Move straight down in Cartesian space until a collision is detected.

    \item Open the gripper.

    \item Rise back to the midpoint waypoint.

    \item Move in joint space to above the place location (slightly higher reference height
          for the return pick).

    \item Move straight down in Cartesian space until a collision is detected.

    \item Close the gripper.

    \item Move in joint space back to above the pick-up location.

    \item Move straight down in Cartesian space until a collision is detected.

    \item Open the gripper --- task complete.
\end{itemize}

Collision detection is used identically for both the descend-to-pick and descend-to-place
phases.  It works by publishing a desired downward Cartesian velocity via inverse kinematics
and monitoring the end-effector $z$-position to determine whether the end-effector is still
moving or has made contact.  The controller sets a downward speed (in Cartesian space), a
low velocity threshold below which motion is deemed to have stopped, and a minimum time for
which the robot must remain stationary before the contact is confirmed.  These two parameters
are necessary to handle the noise and inaccuracies in the reported joint states.

Gripper-force regulation is implemented in a similar fashion.  The gripper is commanded
toward a closed target position; once the gripper velocity falls below a threshold for a
specified hold time, the controller assumes the object is grasped and freezes the gripper
command at its current position.  This prevents the gripper from crushing the object.

\subsection*{Supporting code}

The controller is implemented in \texttt{comp.cpp} and executed via the scripts
\texttt{comp\_hw.sh} and \texttt{comp\_sim.sh}.  The source code is located in
\texttt{/ros\_ws/src/controllers\_new/}.  There are also the files
\texttt{kinematics.hpp/.cpp} which contain all the mathematics required to determine
trajectories.  This includes the Jacobian, inverse velocity kinematics and inverse
kinematics.

The code for the kinematics is structured as follows.  At its foundation there is the
virtual base class \texttt{RobotPart}, from which the \texttt{Link} and \texttt{Joint}
classes are derived.  From the \texttt{Joint} class the \texttt{Revolute} class is derived,
implementing the joint interface.  These objects are used by the generic
\texttt{SequentialRobot} class, which encapsulates all forward and inverse kinematics of
any sequential robot.  In the pick-and-place controller this class is also used for state
tracking: the joint states reported by the \texttt{joint\_states} topic are forwarded to a
\texttt{SequentialRobot} object, which then always has the current forward kinematics and
the Jacobian available.

The inverse kinematics implementation is an iterative method based on the Jacobian.  An
initial guess can be provided by setting the robot state before calling
\texttt{required\_joint\_angles}.  From this initial guess, the state is incrementally
pushed towards the desired pose.  This is done by computing the pose error, which is
interpreted as a desired end-effector velocity pointing towards the target pose.  The
Jacobian is then used to solve for the required joint velocities (via
\texttt{solve\_for\_joint\_velocities\_dls}), which are multiplied by a learning rate and
added to the joint positions (implicitly clamping between the limits in the \texttt{Joint}
class).  This repeats until convergence or divergence; if divergence occurs, the solver
retries with random initial guesses to improve the chance of finding a solution.

The pose error is a 6D vector of the position error and the rotation error.  The position
error is trivial; the rotation error is obtained from the current transform from forward
kinematics and the desired transform from Euler angles.  The rotation parts are converted to
quaternions, where the difference between them is used as the error.  This error quaternion
is converted to an axis-angle representation which represents the rotational velocity towards
the desired pose.

\subsection*{Results}

The robot successfully completed the full pick-and-place cycle with a fragile object.  The
collision-detection approach proved reliable for locating the surface of the platform without
requiring knowledge of its exact height.  The gripper-force regulation prevented excessive
squeezing: once the object resisted further closure the gripper command was frozen, and the
object arrived undamaged at the destination.

The main challenges observed were oscillations in the joint positions when holding the
gripper stationary during the descent phase, and occasional IK failures when the target
pose was near a kinematic singularity.  Performance could be improved by switching to
velocity control mode for the approach phases, which would allow much smoother and faster
motion without the position-control hunting behaviour.

\end{document}
